% GENERATED FILE. Edit canonical lessons, then run npm run generate:books.
% canonical-source-hash: fnv1a64:f22f1f8490f033ab
% canonical-lessons: GU-C179-yojna, GU-C179-nakki-karvu, GU-C179-karyakram, GU-C179-pahelethi, GU-C179-cokkas

\chapter{Plans and programmes}
\label{ch:gu-plans-and-programmes}
\hlchaptermodality{179}

\begin{chapteropening}
\textbf{By the end of this chapter:} \emph{I can say what we have decided and what is planned.}

Complete the last lesson of chapter 179: I can say what we have decided and what is planned.
\end{chapteropening}

\section[\texorpdfstring{\gu{યોજના} (yojnā)}{યોજના (yojnā)}]{\gu{યોજના} (yojnā) — plan, scheme}

\label{lesson:GU-C179-yojna}



\begin{quote}
Before the new one: say the Gujarati for remaining, then the Gujarati for to go by.
\end{quote}

\subsection*{You'll want to know: \gu{યોજના}}
\textbf{\gu{યોજના}} — \emph{yojnā} — \textquotedblleft{}plan, scheme\textquotedblright{}.

\begin{quote}
\textbf{\gu{અમારી} \gu{યોજના} \gu{કાલે} \gu{શરૂ} \gu{થશે}.} — \emph{amārī yojnā kāle śarū thaśe} — \textquotedblleft{}Our plan will start tomorrow.\textquotedblright{}
\end{quote}

\begin{cousinweb}
\textbf{\gu{યોજના}} (\emph{yojnā}) is Sanskrit \emph{yojanā}, \textquotedblleft{}joining, arranging\textquotedblright{}, taken over whole as a learned loan, a tatsama, on the root \emph{yuj}, \textquotedblleft{}to join\textquotedblright{}.
\end{cousinweb}

\begin{grammarlens}[title={What you've built}]
One more word for everyday life at A2.
\end{grammarlens}

\subsection*{Your turn}
\begin{itemize}
\raggedright
  \item \emph{Say it:} \emph{yojnā}
  \item \emph{Say it:} \emph{yojnā}, once more
  \item \emph{Say it:} \emph{amārī yojnā kāle śarū thaśe} — Our plan will start tomorrow.
  \item \emph{Recall:} say the Gujarati for remaining, then the Gujarati for to go by, then say \emph{yojnā} again
\end{itemize}

\begin{tcolorbox}[breakable,colback=teal!4,colframe=teal!35!black,title={Before you move on}]
What does \gu{યોજના} mean? (\textquotedblleft{}Plan, scheme\textquotedblright{}.) Say it once more.
\end{tcolorbox}
\section[\texorpdfstring{\gu{નક્કી} \gu{કરવું} (nakkī karvũ)}{નક્કી કરવું (nakkī karvũ)}]{\gu{નક્કી} \gu{કરવું} (nakkī karvũ) — to decide, to fix}

\label{lesson:GU-C179-nakki-karvu}



\begin{quote}
Before the new one: say the Gujarati for to go by, then the Gujarati for plan.
\end{quote}

\subsection*{You'll want to know: \gu{નક્કી} \gu{કરવું}}
\textbf{\gu{નક્કી} \gu{કરવું}} — \emph{nakkī karvũ} — \textquotedblleft{}to decide, to fix\textquotedblright{}.

\begin{quote}
\textbf{\gu{અમે} \gu{કાલે} \gu{તારીખ} \gu{નક્કી} \gu{કરીશું}.} — \emph{ame kāle tārīkh nakkī karīśũ} — \textquotedblleft{}We will fix the date tomorrow.\textquotedblright{}
\end{quote}

\begin{cousinweb}
\textbf{\gu{નક્કી} \gu{કરવું}} (\emph{nakkī karvũ}) is built from \emph{nakkī}, \textquotedblleft{}settled, fixed\textquotedblright{}, and the inherited \emph{karvũ}, \textquotedblleft{}to do\textquotedblright{}. The origin of \emph{nakkī} is not settled: CDIAL lists no Gujarati form, and EtymDB has no entry.
\end{cousinweb}

\begin{grammarlens}[title={What you've built}]
One more word for everyday life at A2.
\end{grammarlens}

\subsection*{Your turn}
\begin{itemize}
\raggedright
  \item \emph{Say it:} \emph{nakkī karvũ}
  \item \emph{Say it:} \emph{nakkī karvũ}, once more
  \item \emph{Say it:} \emph{ame kāle tārīkh nakkī karīśũ} — We will fix the date tomorrow.
  \item \emph{Recall:} say the Gujarati for to go by, then the Gujarati for plan, then say \emph{nakkī karvũ} again
\end{itemize}

\begin{tcolorbox}[breakable,colback=teal!4,colframe=teal!35!black,title={Before you move on}]
What does \gu{નક્કી} \gu{કરવું} mean? (\textquotedblleft{}To decide, to fix\textquotedblright{}.) Say it once more.
\end{tcolorbox}
\section[\texorpdfstring{\gu{કાર્યક્રમ} (kāryakram)}{કાર્યક્રમ (kāryakram)}]{\gu{કાર્યક્રમ} (kāryakram) — programme, event}

\label{lesson:GU-C179-karyakram}



\begin{quote}
Before the new one: say the Gujarati for plan, then the Gujarati for to decide.
\end{quote}

\subsection*{You'll want to know: \gu{કાર્યક્રમ}}
\textbf{\gu{કાર્યક્રમ}} — \emph{kāryakram} — \textquotedblleft{}programme, event\textquotedblright{}.

\begin{quote}
\textbf{\gu{શાળામાં} \gu{સાંજે} \gu{કાર્યક્રમ} \gu{હશે}.} — \emph{śāḷāmā̃ sā̃je kāryakram haśe} — \textquotedblleft{}There will be a programme at the school in the evening.\textquotedblright{}
\end{quote}

\begin{cousinweb}
\textbf{\gu{કાર્યક્રમ}} (\emph{kāryakram}) is a learned compound built from Sanskrit \emph{kārya}, \textquotedblleft{}work, task\textquotedblright{}, and \emph{krama}, \textquotedblleft{}step, order\textquotedblright{}, a tatsama.
\end{cousinweb}

\begin{grammarlens}[title={What you've built}]
One more word for everyday life at A2.
\end{grammarlens}

\subsection*{Your turn}
\begin{itemize}
\raggedright
  \item \emph{Say it:} \emph{kāryakram}
  \item \emph{Say it:} \emph{kāryakram}, once more
  \item \emph{Say it:} \emph{śāḷāmā̃ sā̃je kāryakram haśe} — There will be a programme at the school in the evening.
  \item \emph{Recall:} say the Gujarati for plan, then the Gujarati for to decide, then say \emph{kāryakram} again
\end{itemize}

\begin{tcolorbox}[breakable,colback=teal!4,colframe=teal!35!black,title={Before you move on}]
What does \gu{કાર્યક્રમ} mean? (\textquotedblleft{}Programme, event\textquotedblright{}.) Say it once more.
\end{tcolorbox}
\section[\texorpdfstring{\gu{પહેલેથી} (pahelethī)}{પહેલેથી (pahelethī)}]{\gu{પહેલેથી} (pahelethī) — beforehand, in advance}

\label{lesson:GU-C179-pahelethi}



\begin{quote}
Before the new one: say the Gujarati for to decide, then the Gujarati for programme.
\end{quote}

\subsection*{You'll want to know: \gu{પહેલેથી}}
\textbf{\gu{પહેલેથી}} — \emph{pahelethī} — \textquotedblleft{}beforehand, in advance\textquotedblright{}.

\begin{quote}
\textbf{\gu{હું} \gu{પહેલેથી} \gu{ટિકિટ} \gu{લઈશ}.} — \emph{hũ pahelethī ṭikiṭ laīś} — \textquotedblleft{}I will get the ticket in advance.\textquotedblright{}
\end{quote}

\begin{cousinweb}
\textbf{\gu{પહેલેથી}} (\emph{pahelethī}) is built on the inherited \emph{pahelũ}, \textquotedblleft{}first\textquotedblright{}, with \emph{-thī}, \textquotedblleft{}from\textquotedblright{}, so it means from the first.
\end{cousinweb}

\begin{grammarlens}[title={What you've built}]
One more word for everyday life at A2.
\end{grammarlens}

\subsection*{Your turn}
\begin{itemize}
\raggedright
  \item \emph{Say it:} \emph{pahelethī}
  \item \emph{Say it:} \emph{pahelethī}, once more
  \item \emph{Say it:} \emph{hũ pahelethī ṭikiṭ laīś} — I will get the ticket in advance.
  \item \emph{Recall:} say the Gujarati for to decide, then the Gujarati for programme, then say \emph{pahelethī} again
\end{itemize}

\begin{tcolorbox}[breakable,colback=teal!4,colframe=teal!35!black,title={Before you move on}]
What does \gu{પહેલેથી} mean? (\textquotedblleft{}Beforehand, in advance\textquotedblright{}.) Say it once more.
\end{tcolorbox}
\section[\texorpdfstring{\gu{ચોક્કસ} (cokkas)}{ચોક્કસ (cokkas)}]{\gu{ચોક્કસ} (cokkas) — certain, sure; exact}

\label{lesson:GU-C179-cokkas}



\begin{quote}
Before the new one: say the Gujarati for programme, then the Gujarati for beforehand.
\end{quote}

\subsection*{You'll want to know: \gu{ચોક્કસ}}
\textbf{\gu{ચોક્કસ}} — \emph{cokkas} — \textquotedblleft{}certain, sure; exact\textquotedblright{}.

\begin{quote}
\textbf{\gu{હું} \gu{ચોક્કસ} \gu{આવીશ}.} — \emph{hũ cokkas āvīś} — \textquotedblleft{}I will certainly come.\textquotedblright{}
\end{quote}

\begin{cousinweb}
The origin of \textbf{\gu{ચોક્કસ}} (\emph{cokkas}) is not settled. CDIAL lists no Gujarati form of it, and EtymDB has no entry for it.
\end{cousinweb}

\begin{grammarlens}[title={What you've built}]
One more word for everyday life at A2.
\end{grammarlens}

\subsection*{Your turn}
\begin{itemize}
\raggedright
  \item \emph{Say it:} \emph{cokkas}
  \item \emph{Say it:} \emph{cokkas}, once more
  \item \emph{Say it:} \emph{hũ cokkas āvīś} — I will certainly come.
  \item \emph{Recall:} say the Gujarati for programme, then the Gujarati for beforehand, then say \emph{cokkas} again
\end{itemize}

\begin{tcolorbox}[breakable,colback=teal!4,colframe=teal!35!black,title={Before you move on}]
What does \gu{ચોક્કસ} mean? (\textquotedblleft{}Certain, sure; exact\textquotedblright{}.) Say it once more.
\end{tcolorbox}
